\section{模型评价与推广}

\subsection{模型优点}

\begin{enumerate}
  \item \textbf{结构性解释}：20 点定案布局的每类点都有明确的物理/几何依据——7 覆盖基点
        保全向覆盖、12 均匀外圈点射向任意贴边与内部半径源的迎光区、原点零成本起点扫描；
        全流程（检测—定位—清除）参数全部集中一处（\texttt{cumcm/t4/config.py}），
        "一处改、全局生效"。
  \item \textbf{统计验证充分}：$422.5$ 万例（真实随机 $400$ 万 $+$ 贴边对抗 $4.3$ 万
        $+$ 精细对抗 $18$ 万）三重口径听到率 $99.9891\%$，贴边对抗 0 漏；20 局演练
        $256/256$ 全清，定位误差全部落在清除半径内。
  \item \textbf{确定性与可复现}：所有启发式（最近邻 + 2-opt、补测选点、顺路清除扇区）
        无随机算子，同 seed 逐字节可复现，A/B 对拍（时间旋钮、布局对比）结论可信。
  \item \textbf{时间最优性有独立证据}：6950 s 级的下界不仅靠 20 局裁决，还由完全独立于
        领域先验的神经网络搜索交叉验证（全部候选更慢），排除了"人工遗漏最优解"的疑虑。
\end{enumerate}

\subsection{模型缺点}

\begin{enumerate}
  \item \textbf{非严格保证}：由于 $g,\theta,R$ 连续未知，检测保证为实测统计性
        （残余漏例约 $0.01\%$，集中在"内部半径源 + 朝外波束 + 方向恰落 7 基点方位空隙"），
        未做到数学上的严格不漏（覆盖率下界仅对全向源成立）。
  \item \textbf{依赖 7 基点结构}：布局复用问题三的 7 覆盖基点，若题目改为任意多边形区域或
        源生成域变化，基点结构与外圈半径 $1850$ m 需重新标定。
  \item \textbf{神经网络失效的统计墙}：$10$ 万级蒙特卡洛精评无法分辨 $0.005\%$ 级听率差异，
        神经网络在"极低漏率精调"上受限于采样分辨率。
\end{enumerate}

\subsection{推广}

\begin{enumerate}
  \item \textbf{K 圆覆盖清除}（定位区域用至多 $4$ 个半径 $20$ m 圆覆盖逐点试清）可推广到
        任意"区域约束 + 清除半径"的搜索清除任务（如无人机搜索、矿难定位）。
  \item \textbf{顺路清除机制}（站间方位扇区尝试）适用于一切"巡视 + 定点作业"的路径—任务
        联合优化场景，扇区半张角按任务成功率与绕路代价实测标定。
  \item \textbf{迎光区覆盖思想}：对"扇形视域 + 未知朝向"的目标，把测量点放在目标可能出现
        区域的\textbf{视域纵深内沿}（本文 $1850$ m $\leftarrow$ 贴边内沿 $1770$ m），是
        通用且可迁移的几何结论。
  \item 完整求解链路（\texttt{T1.py}/\texttt{T3.py}/\texttt{T4.py} + 本地模拟器回放 + 同
        seed 对拍）可服务任意"官方模拟器受限"的竞赛场景，作为离线验证与证据链基础设施。
\end{enumerate}